% ID: SEL-D0074-OB-E03
% ESTADO: cotejada. Segunda pasada visual y borrador compilado revisados.
% Fuente: Selectividad/sel_2008_matematicas_aplicadas/Mat CCSS 2 - Examen y criterios de correción.pdf
% SHA-256: 6439a5956271f683dc45ba870695a7399db4412c9b0eaf1a35f908a3857ae648
% Página: 2; opción B; ejercicio original 3.

\noindent\underline{Parte I}

Se consideran los sucesos $A$ y $B$.
\begin{enumerate}[label=\alph*)]
\item \textbf{(0.75 puntos)} Exprese, utilizando las operaciones con sucesos, los siguientes sucesos:
\begin{enumerate}[label=\arabic*.]
\item Que no ocurra ninguno de los dos.
\item Que ocurra al menos uno de los dos.
\item Que ocurra $B$, pero que no ocurra $A$.
\end{enumerate}
\item \textbf{(1.25 puntos)} Sabiendo que $P(A)=0.5$, $P(B)=0.5$ y $P(A/B)=0.3$, halle $P(A\cup B)$.
\end{enumerate}
\par\noindent\begin{minipage}{\linewidth}
\noindent\underline{Parte II}

\textbf{(2 puntos)} Se ha aplicado un medicamento a una muestra de 200 enfermos y se ha observado una respuesta positiva en 140 de ellos. Estímese, mediante un intervalo de confianza del $99\%$, la proporción de enfermos que responderían positivamente si este medicamento se aplicase a la población de la que se ha extraído la muestra.
\end{minipage}\par
